\documentclass[12pt,a4paper]{article}

% Greek language support
\usepackage[utf8]{inputenc}
\usepackage[LGR,T1]{fontenc}
\usepackage[greek,english]{babel}
\usepackage{alphabeta}

% Packages
\usepackage{geometry}
\usepackage{graphicx}
\usepackage{booktabs}
\usepackage{longtable}
\usepackage{array}
\usepackage{multirow}
\usepackage{float}
\usepackage{caption}
\usepackage{subcaption}
\usepackage{xcolor}
\usepackage{listings}
\usepackage{amsmath}
\usepackage{amssymb}
\usepackage{hyperref}
\usepackage{fancyhdr}
\usepackage{titlesec}
\usepackage{enumitem}
\usepackage{tcolorbox}

% Page settings
\geometry{
    left=2.5cm,
    right=2.5cm,
    top=2.5cm,
    bottom=2.5cm
}

% Python code settings
\definecolor{codegreen}{rgb}{0,0.6,0}
\definecolor{codegray}{rgb}{0.5,0.5,0.5}
\definecolor{codepurple}{rgb}{0.58,0,0.82}
\definecolor{backcolour}{rgb}{0.95,0.95,0.92}

\lstdefinestyle{pythonstyle}{
    backgroundcolor=\color{backcolour},
    commentstyle=\color{codegreen},
    keywordstyle=\color{blue},
    numberstyle=\tiny\color{codegray},
    stringstyle=\color{codepurple},
    basicstyle=\ttfamily\footnotesize,
    breakatwhitespace=false,
    breaklines=true,
    captionpos=b,
    keepspaces=true,
    numbers=left,
    numbersep=5pt,
    showspaces=false,
    showstringspaces=false,
    showtabs=false,
    tabsize=2,
    language=Python,
    frame=single,
    rulecolor=\color{black},
    extendedchars=false,
    literate={<=}{{$\leq$}}1 {>=}{{$\geq$}}1
}

\lstset{style=pythonstyle}

% Hyperref settings
\hypersetup{
    colorlinks=true,
    linkcolor=blue,
    filecolor=magenta,
    urlcolor=cyan,
    citecolor=green
}

% Header/footer settings
\pagestyle{fancy}
\fancyhf{}
\rhead{AEM: 6931}
\lhead{Data Analysis 2025-26}
\cfoot{\thepage}

% Custom boxes
\newtcolorbox{infobox}[1][]{
    colback=blue!5!white,
    colframe=blue!75!black,
    fonttitle=\bfseries,
    title=#1
}

\newtcolorbox{warningbox}[1][]{
    colback=orange!5!white,
    colframe=orange!75!black,
    fonttitle=\bfseries,
    title=#1
}

\newtcolorbox{successbox}[1][]{
    colback=green!5!white,
    colframe=green!75!black,
    fonttitle=\bfseries,
    title=#1
}

% Title
\title{
    \vspace{-1cm}
    \textbf{\textgreek{Ανάλυση Δεδομένων} 2025-26}\\
    \large \textgreek{1η Εργασία: Καθαρισμός και Ανάλυση Δεδομένων Ποιότητας Κρασιού}\\
    \vspace{0.5cm}
    \normalsize \textgreek{Ερώτημα 1: Δειγματοληψία και Καθαρισμός Δεδομένων}
}

\author{
    \textbf{AEM:} 6931
}

\date{\textgreek{Μάρτιος} 2026}

\begin{document}

\maketitle

\tableofcontents
\newpage

% ==============================================================================
\section{\textgreek{Εισαγωγή}}
% ==============================================================================

\subsection{\textgreek{Σκοπός της Εργασίας}}

\textgreek{Η παρούσα εργασία αφορά την ανάλυση ενός συνόλου δεδομένων (dataset) που περιέχει χαρακτηριστικά ποιότητας διαφόρων τύπων κρασιού. Ο στόχος του πρώτου ερωτήματος είναι:}

\begin{enumerate}
    \item \textgreek{Η επιλογή ενός} \textbf{\textgreek{τυχαίου δείγματος}} \textgreek{που αποτελεί το 80\% των αρχικών δεδομένων}
    \item \textgreek{Ο} \textbf{\textgreek{έλεγχος ποιότητας}} \textgreek{των δεδομένων για τον εντοπισμό προβλημάτων}
    \item \textgreek{Ο} \textbf{\textgreek{καθαρισμός}} \textgreek{των δεδομένων με κατάλληλες τεχνικές}
\end{enumerate}

\subsection{\textgreek{Τι είναι ο Καθαρισμός Δεδομένων;}}

\begin{infobox}[\textgreek{Ορισμός: Καθαρισμός Δεδομένων (Data Cleaning)}]
\textgreek{Ο καθαρισμός δεδομένων είναι η διαδικασία εντοπισμού και διόρθωσης (ή αφαίρεσης) λανθασμένων, ελλιπών, διπλότυπων ή ασυνεπών εγγραφών από ένα σύνολο δεδομένων. Είναι ένα κρίσιμο βήμα πριν από οποιαδήποτε ανάλυση, καθώς η ποιότητα των αποτελεσμάτων εξαρτάται άμεσα από την ποιότητα των δεδομένων.}
\end{infobox}

\subsection{\textgreek{Γιατί είναι Σημαντικός ο Καθαρισμός Δεδομένων;}}

\textgreek{Τα πραγματικά δεδομένα σπάνια είναι τέλεια. Συνήθη προβλήματα περιλαμβάνουν:}

\begin{itemize}
    \item \textbf{\textgreek{Τυπογραφικά λάθη:}} \textgreek{Λανθασμένη πληκτρολόγηση (π.χ. ``whte'' αντί ``white'')}
    \item \textbf{\textgreek{Ελλείπουσες τιμές (Missing Values):}} \textgreek{Κενά πεδία σε ορισμένες εγγραφές}
    \item \textbf{\textgreek{Μη έγκυρες τιμές:}} \textgreek{Τιμές εκτός λογικών ορίων (π.χ. pH = 999)}
    \item \textbf{\textgreek{Διπλοεγγραφές:}} \textgreek{Η ίδια εγγραφή εμφανίζεται πολλές φορές}
\end{itemize}

\textgreek{Αν δεν αντιμετωπιστούν αυτά τα προβλήματα, τα αποτελέσματα της ανάλυσης θα είναι αναξιόπιστα ή παραπλανητικά.}

% ==============================================================================
\section{\textgreek{Περιγραφή του Συνόλου Δεδομένων}}
% ==============================================================================

\subsection{\textgreek{Πηγή και Περιεχόμενο}}

\textgreek{Το σύνολο δεδομένων περιέχει μετρήσεις χημικών χαρακτηριστικών κρασιών (κόκκινων και λευκών) μαζί με βαθμολογίες ποιότητας. Κάθε γραμμή αντιπροσωπεύει ένα δείγμα κρασιού.}

\subsection{\textgreek{Μεταβλητές του Dataset}}

\begin{table}[H]
\centering
\caption{\textgreek{Περιγραφή των μεταβλητών του dataset}}
\label{tab:variables}
\begin{tabular}{@{}clp{7cm}@{}}
\toprule
\textbf{A/A} & \textbf{\textgreek{Μεταβλητή}} & \textbf{\textgreek{Περιγραφή}} \\
\midrule
1 & Fixed acidity & \textgreek{Σταθερή οξύτητα - Μη πτητικά οξέα (g/L)} \\
2 & Volatile acidity & \textgreek{Πτητική οξύτητα - Οξικό οξύ (g/L)} \\
3 & Citric acid & \textgreek{Κιτρικό οξύ (g/L)} \\
4 & Residual sugar & \textgreek{Υπολειμματικά σάκχαρα μετά τη ζύμωση (g/L)} \\
5 & Chlorides & \textgreek{Χλωριούχα άλατα - Δείκτης αλατότητας (g/L)} \\
6 & Free sulfur dioxide & \textgreek{Ελεύθερο SO\textsubscript{2} (mg/L)} \\
7 & Total sulfur dioxide & \textgreek{Συνολικό SO\textsubscript{2} (mg/L)} \\
8 & Density & \textgreek{Πυκνότητα (g/cm\textsuperscript{3})} \\
9 & pH & \textgreek{Μέτρο της οξύτητας (κλίμακα 0-14)} \\
10 & Sulphates & \textgreek{Θειικά άλατα (g/L)} \\
11 & Alcohol & \textgreek{Περιεκτικότητα σε αιθανόλη (\% vol)} \\
12 & Quality & \textgreek{Βαθμολογία ποιότητας (κλίμακα 0-10)} \\
13 & Wine\_type & \textgreek{Τύπος κρασιού: red ή white} \\
\bottomrule
\end{tabular}
\end{table}

\begin{infobox}[\textgreek{Τύποι Μεταβλητών}]
\begin{itemize}
    \item \textbf{\textgreek{Συνεχείς (Continuous):}} \textgreek{Μεταβλητές 1-11 - μπορούν να πάρουν οποιαδήποτε αριθμητική τιμή σε ένα εύρος}
    \item \textbf{\textgreek{Διακριτές/Κατηγορικές:}} \textgreek{Quality (διακριτή κλίμακα 0-10), Wine\_type (κατηγορική)}
\end{itemize}
\end{infobox}

\subsection{\textgreek{Αρχικές Διαστάσεις}}

\textgreek{Το αρχικό dataset περιείχε:}
\begin{itemize}
    \item \textbf{8.385 \textgreek{γραμμές}} \textgreek{(εγγραφές/δείγματα κρασιού)}
    \item \textbf{15 \textgreek{στήλες}} \textgreek{(13 μεταβλητές + 2 κενές στήλες)}
\end{itemize}

% ==============================================================================
\section{\textgreek{Δειγματοληψία (Sampling)}}
% ==============================================================================

\subsection{\textgreek{Τι είναι η Δειγματοληψία;}}

\begin{infobox}[\textgreek{Ορισμός: Δειγματοληψία}]
\textgreek{Η δειγματοληψία είναι η διαδικασία επιλογής ενός υποσυνόλου (δείγματος) από ένα μεγαλύτερο σύνολο δεδομένων (πληθυσμό). Χρησιμοποιείται για να μειώσει τον όγκο των δεδομένων ή για να δημιουργήσει διαφορετικά σύνολα για εκπαίδευση και αξιολόγηση μοντέλων.}
\end{infobox}

\subsection{\textgreek{Γιατί τυχαία δειγματοληψία;}}

\textgreek{Η} \textbf{\textgreek{τυχαία δειγματοληψία}} \textgreek{εξασφαλίζει ότι:}
\begin{enumerate}
    \item \textgreek{Κάθε εγγραφή έχει} \textbf{\textgreek{ίση πιθανότητα}} \textgreek{να επιλεγεί}
    \item \textgreek{Το δείγμα είναι} \textbf{\textgreek{αντιπροσωπευτικό}} \textgreek{του συνόλου}
    \item \textgreek{Αποφεύγεται η} \textbf{\textgreek{μεροληψία}} (bias) \textgreek{στην επιλογή}
\end{enumerate}

\subsection{\textgreek{Ο Ρόλος του Seed (Σπόρος)}}

\begin{warningbox}[\textgreek{Σημαντικό: Random Seed}]
\textgreek{Οι υπολογιστές δημιουργούν ``ψευδοτυχαίους'' αριθμούς χρησιμοποιώντας αλγορίθμους. Ο} \textbf{seed} \textgreek{(σπόρος) είναι μια αρχική τιμή που καθορίζει την ακολουθία των τυχαίων αριθμών.}

\textbf{\textgreek{Γιατί είναι σημαντικό:}} \textgreek{Χρησιμοποιώντας τον ίδιο seed, παίρνουμε πάντα τα ίδια ``τυχαία'' αποτελέσματα. Αυτό επιτρέπει την} \textbf{\textgreek{αναπαραγωγιμότητα}} (reproducibility) \textgreek{της ανάλυσης.}

\textgreek{Στην εργασία, χρησιμοποιούμε ως seed τον ΑΕΜ (6931), ώστε κάθε φοιτητής να εργάζεται με διαφορετικό δείγμα.}
\end{warningbox}

\subsection{\textgreek{Υλοποίηση στην Python}}

\begin{lstlisting}[caption={Sampling code}]
import pandas as pd

# Load data
df_original = pd.read_csv('Data Analysis_2026 1st Case_Data.csv')

# Define parameters
AEM = 6931              # The seed
SAMPLE_FRACTION = 0.80  # Sample percentage (80%)

# Sampling
df = df_original.sample(frac=SAMPLE_FRACTION, random_state=AEM)

# Reset index
df = df.reset_index(drop=True)
\end{lstlisting}

\textbf{\textgreek{Επεξήγηση του κώδικα:}}
\begin{itemize}
    \item \texttt{sample()}: \textgreek{Η συνάρτηση της pandas για τυχαία επιλογή γραμμών}
    \item \texttt{frac=0.80}: \textgreek{Επιλέγουμε το 80\% των γραμμών}
    \item \texttt{random\_state=AEM}: \textgreek{Ορίζουμε τον seed για αναπαραγωγιμότητα}
    \item \texttt{reset\_index()}: \textgreek{Επαναριθμούμε τις γραμμές από 0}
\end{itemize}

\subsection{\textgreek{Αποτέλεσμα Δειγματοληψίας}}

\begin{table}[H]
\centering
\caption{\textgreek{Αποτελέσματα δειγματοληψίας}}
\begin{tabular}{@{}lcc@{}}
\toprule
\textbf{\textgreek{Μέγεθος}} & \textbf{\textgreek{Γραμμές}} & \textbf{\textgreek{Ποσοστό}} \\
\midrule
\textgreek{Αρχικό dataset} & 8.385 & 100\% \\
\textgreek{Δείγμα (ΑΕΜ: 6931)} & 6.708 & 80\% \\
\bottomrule
\end{tabular}
\end{table}

% ==============================================================================
\section{\textgreek{Εντοπισμός Προβλημάτων στα Δεδομένα}}
% ==============================================================================

\textgreek{Μετά τη δειγματοληψία, πραγματοποιήθηκε συστηματικός έλεγχος για τον εντοπισμό προβλημάτων.}

\subsection{\textgreek{Κατηγορία 1: Τυπογραφικά Λάθη στη Στήλη wine\_type}}

\subsubsection{\textgreek{Τι ελέγξαμε}}

\textgreek{Η στήλη} \texttt{wine\_type} \textgreek{πρέπει να περιέχει μόνο δύο τιμές: ``red'' ή ``white''.}

\begin{lstlisting}[caption={Check unique values}]
# Display all unique values
print(df['wine_type'].unique())
\end{lstlisting}

\subsubsection{\textgreek{Τι βρήκαμε}}

\begin{table}[H]
\centering
\caption{\textgreek{Τυπογραφικά λάθη στη στήλη wine\_type}}
\label{tab:typos}
\begin{tabular}{@{}lcc@{}}
\toprule
\textbf{\textgreek{Λανθασμένη Τιμή}} & \textbf{\textgreek{Πλήθος}} & \textbf{\textgreek{Σωστή Τιμή}} \\
\midrule
Redd  & 75 & red \\
whit & 90 & white \\
r3d & 64 & red \\
WHITTE & 77 & white \\
whate & 39 & white \\
999 & 59 & \textgreek{Μη έγκυρη} \\
9999 & 45 & \textgreek{Μη έγκυρη} \\
-999 & 64 & \textgreek{Μη έγκυρη} \\
? & 47 & \textgreek{Μη έγκυρη} \\
\midrule
\textbf{\textgreek{Σύνολο}} & \textbf{560} & \\
\bottomrule
\end{tabular}
\end{table}

\begin{infobox}[\textgreek{Ανάλυση των λαθών}]
\begin{itemize}
    \item \textbf{``Redd '', ``r3d''}: \textgreek{Παραλλαγές του ``red'' με τυπογραφικά λάθη}
    \item \textbf{``whit'', ``WHITTE'', ``whate''}: \textgreek{Παραλλαγές του ``white''}
    \item \textbf{``999'', ``9999'', ``-999'', ``?''}: \textgreek{Κωδικοί για ελλείπουσες τιμές}
\end{itemize}
\end{infobox}

\subsection{\textgreek{Κατηγορία 2: Μη Αριθμητικές Τιμές}}

\subsubsection{\textgreek{Τι ελέγξαμε}}

\textgreek{Οι 12 μεταβλητές (εκτός του wine\_type) πρέπει να περιέχουν αριθμητικές τιμές.}

\begin{lstlisting}[caption={Convert to numeric}]
# Convert with errors='coerce' converts non-numeric to NaN
df['fixed acidity'] = pd.to_numeric(df['fixed acidity'], errors='coerce')
\end{lstlisting}

\subsubsection{\textgreek{Τι βρήκαμε}}

\begin{table}[H]
\centering
\caption{\textgreek{Μη αριθμητικές τιμές ανά στήλη}}
\label{tab:non_numeric}
\begin{tabular}{@{}lc@{}}
\toprule
\textbf{\textgreek{Μεταβλητή}} & \textbf{\textgreek{Μη αριθμητικές τιμές}} \\
\midrule
fixed acidity & 87 \\
volatile acidity & 45 \\
citric acid & 61 \\
residual sugar & 56 \\
chlorides & 57 \\
free sulfur dioxide & 56 \\
total sulfur dioxide & 53 \\
density & 53 \\
pH & 151 \\
sulphates & 63 \\
alcohol & 105 \\
quality & 58 \\
\midrule
\textbf{\textgreek{Σύνολο}} & \textbf{845} \\
\bottomrule
\end{tabular}
\end{table}

\subsection{\textgreek{Κατηγορία 3: Αδύνατες/Ακραίες Τιμές}}

\subsubsection{\textgreek{Τι ελέγξαμε}}

\textgreek{Ορίσαμε λογικά όρια για κάθε μεταβλητή βάσει επιστημονικής γνώσης:}

\begin{table}[H]
\centering
\caption{\textgreek{Λογικά όρια τιμών ανά μεταβλητή}}
\label{tab:limits}
\begin{tabular}{@{}lccp{4.5cm}@{}}
\toprule
\textbf{\textgreek{Μεταβλητή}} & \textbf{Min} & \textbf{Max} & \textbf{\textgreek{Αιτιολόγηση}} \\
\midrule
Fixed acidity & 0 & 20 & \textgreek{Τυπικές τιμές 4-16 g/L} \\
Volatile acidity & 0 & 5 & \textgreek{Τυπικές τιμές 0.1-1.5 g/L} \\
Citric acid & 0 & 3 & \textgreek{Τυπικές τιμές 0-1 g/L} \\
Residual sugar & 0 & 100 & \textgreek{Ξηρά <4, γλυκά ως 100 g/L} \\
Chlorides & 0 & 1 & \textgreek{Τυπικές τιμές 0.01-0.2 g/L} \\
Free SO\textsubscript{2} & 0 & 300 & \textgreek{Νομικό όριο ~300 mg/L} \\
Total SO\textsubscript{2} & 0 & 500 & \textgreek{Νομικό όριο ~400 mg/L} \\
Density & 0.9 & 1.1 & \textgreek{Κοντά στο νερό (~1 g/cm³)} \\
pH & 2 & 5 & \textgreek{Κρασί: τυπικά 2.9-4.0} \\
Sulphates & 0 & 3 & \textgreek{Τυπικές τιμές 0.3-1.5 g/L} \\
Alcohol & 5 & 20 & \textgreek{Κρασί: τυπικά 8-15\%} \\
Quality & 0 & 10 & \textgreek{Κλίμακα βαθμολόγησης} \\
\bottomrule
\end{tabular}
\end{table}

\subsubsection{\textgreek{Τι βρήκαμε}}

\begin{table}[H]
\centering
\caption{\textgreek{Αδύνατες τιμές ανά μεταβλητή}}
\label{tab:impossible}
\begin{tabular}{@{}lcp{5cm}@{}}
\toprule
\textbf{\textgreek{Μεταβλητή}} & \textbf{\textgreek{Πλήθος}} & \textbf{\textgreek{Παραδείγματα}} \\
\midrule
fixed acidity & 156 & 999, -999, 9999 \\
volatile acidity & 170 & 999, -999, 9999 \\
citric acid & 164 & 999, -999, 9999 \\
residual sugar & 165 & 999, -999, 9999 \\
chlorides & 161 & 999, -999, 9999 \\
free sulfur dioxide & 151 & 999, -999, 9999 \\
total sulfur dioxide & 160 & 999, -999, 9999 \\
density & 154 & 999, -999, 9999 \\
pH & 202 & 0, 999, -999, 9999 \\
sulphates & 160 & 999, -999, 9999 \\
alcohol & 197 & 999, -999, 9999 \\
quality & 197 & 999, -999, 9999 \\
\midrule
\textbf{\textgreek{Σύνολο}} & \textbf{2.037} & \\
\bottomrule
\end{tabular}
\end{table}

\begin{warningbox}[\textgreek{Παρατήρηση}]
\textgreek{Οι τιμές 999, -999 και 9999 εμφανίζονται συστηματικά σε όλες τις στήλες. Αυτό υποδηλώνει ότι χρησιμοποιήθηκαν} \textbf{\textgreek{ως κωδικοί για ελλείπουσες τιμές}} \textgreek{κατά την καταχώρηση.}
\end{warningbox}

\subsection{\textgreek{Κατηγορία 4: Διπλοεγγραφές (Duplicates)}}

\subsubsection{\textgreek{Τι είναι οι διπλοεγγραφές}}

\textgreek{Μια διπλοεγγραφή είναι μια γραμμή που εμφανίζεται} \textbf{\textgreek{πανομοιότυπη}} \textgreek{με άλλη γραμμή στο dataset.}

\subsubsection{\textgreek{Τι βρήκαμε}}

\begin{lstlisting}[caption={Find duplicates}]
# Find duplicates
duplicates = df.duplicated()
print(f"Number of duplicates: {duplicates.sum()}")
\end{lstlisting}

\begin{successbox}[\textgreek{Αποτέλεσμα}]
\textgreek{Εντοπίστηκαν} \textbf{1.434 \textgreek{διπλοεγγραφές}} \textgreek{στο δείγμα.}
\end{successbox}

\subsection{\textgreek{Κατηγορία 5: Ελλείπουσες Τιμές (Missing Values)}}

\subsubsection{\textgreek{Τι είναι οι ελλείπουσες τιμές}}

\begin{infobox}[\textgreek{Ορισμός: Missing Values}]
\textgreek{Οι ελλείπουσες τιμές (ή NaN - Not a Number) είναι θέσεις στο dataset όπου δεν υπάρχει καταχωρημένη τιμή.}
\end{infobox}

\subsubsection{\textgreek{Τι βρήκαμε}}

\begin{table}[H]
\centering
\caption{\textgreek{Ελλείπουσες τιμές ανά μεταβλητή}}
\label{tab:missing}
\begin{tabular}{@{}lcc@{}}
\toprule
\textbf{\textgreek{Μεταβλητή}} & \textbf{Missing Values} & \textbf{\textgreek{Ποσοστό}} \\
\midrule
fixed acidity & 574 & 10.88\% \\
volatile acidity & 542 & 10.28\% \\
citric acid & 1 & 0.02\% \\
residual sugar & 134 & 2.54\% \\
chlorides & 8 & 0.15\% \\
free sulfur dioxide & 1 & 0.02\% \\
total sulfur dioxide & 1 & 0.02\% \\
density & 545 & 10.33\% \\
pH & 644 & 12.21\% \\
sulphates & 1 & 0.02\% \\
alcohol & 600 & 11.38\% \\
quality & 40 & 0.76\% \\
wine\_type & 1 & 0.02\% \\
\midrule
\textbf{\textgreek{Σύνολο}} & \textbf{3.092} & \\
\bottomrule
\end{tabular}
\end{table}

% ==============================================================================
\section{\textgreek{Αντιμετώπιση Προβλημάτων}}
% ==============================================================================

\subsection{\textgreek{Στρατηγική Αντιμετώπισης}}

\begin{table}[H]
\centering
\caption{\textgreek{Στρατηγική αντιμετώπισης ανά τύπο προβλήματος}}
\label{tab:strategy}
\begin{tabular}{@{}lp{8cm}@{}}
\toprule
\textbf{\textgreek{Πρόβλημα}} & \textbf{\textgreek{Στρατηγική Αντιμετώπισης}} \\
\midrule
\textgreek{Τυπογραφικά λάθη} & \textgreek{Διόρθωση στη σωστή τιμή ή μετατροπή σε NaN} \\
\textgreek{Μη αριθμητικές τιμές} & \textgreek{Μετατροπή σε NaN} \\
\textgreek{Αδύνατες τιμές} & \textgreek{Μετατροπή σε NaN} \\
\textgreek{Διπλοεγγραφές} & \textgreek{Αφαίρεση (κράτηση πρώτης εμφάνισης)} \\
Missing values & \textgreek{Συμπλήρωση με διάμεσο (median) ανά τύπο κρασιού} \\
\bottomrule
\end{tabular}
\end{table}

\subsection{\textgreek{Διόρθωση Τυπογραφικών Λαθών}}

\begin{lstlisting}[caption={Fix typos}]
# Define correction mappings
wine_type_corrections = {
    'Redd ': 'red',     # Typo
    'r3d': 'red',       # Typo (number 3 instead of e)
    'WHITTE': 'white',  # Typo
    'whit': 'white',    # Incomplete word
    'whate': 'white',   # Typo
    '999': np.nan,      # Missing value code
    '9999': np.nan,     # Missing value code
    '-999': np.nan,     # Missing value code
    '?': np.nan,        # Unknown symbol
}

# Apply corrections
for wrong, correct in wine_type_corrections.items():
    mask = df['wine_type'] == wrong
    df.loc[mask, 'wine_type'] = correct
\end{lstlisting}

\begin{successbox}[\textgreek{Αποτέλεσμα}]
\textgreek{Διορθώθηκαν} \textbf{560 \textgreek{εγγραφές}} \textgreek{με τυπογραφικά λάθη.}
\end{successbox}

\subsection{\textgreek{Αντιμετώπιση Αδύνατων Τιμών}}

\begin{lstlisting}[caption={Replace impossible values}]
# Define limits for each variable
value_limits = {
    'fixed acidity': (0, 20),
    'volatile acidity': (0, 5),
    'pH': (2, 5),
    # ... etc
}

# For each variable
for col, (min_val, max_val) in value_limits.items():
    # Find values outside limits
    out_of_range = (df[col] < min_val) | (df[col] > max_val)
    # Replace with NaN
    df.loc[out_of_range, col] = np.nan
\end{lstlisting}

\subsection{\textgreek{Αφαίρεση Διπλοεγγραφών}}

\begin{lstlisting}[caption={Remove duplicates}]
# Remove duplicates, keep first occurrence
df = df.drop_duplicates()
df = df.reset_index(drop=True)
\end{lstlisting}

\begin{successbox}[\textgreek{Αποτέλεσμα}]
\textgreek{Αφαιρέθηκαν} \textbf{1.434 \textgreek{διπλοεγγραφές}}.
\end{successbox}

\subsection{\textgreek{Συμπλήρωση Ελλειπουσών Τιμών}}

\subsubsection{\textgreek{Επιλογή Μεθόδου: Διάμεσος (Median)}}

\begin{infobox}[\textgreek{Γιατί Διάμεσος και όχι Μέσος Όρος;}]
\begin{itemize}
    \item \textgreek{Η} \textbf{\textgreek{διάμεσος}} \textgreek{είναι η μεσαία τιμή όταν τα δεδομένα ταξινομηθούν}
    \item \textgreek{Είναι} \textbf{\textgreek{ανθεκτική}} \textgreek{σε ακραίες τιμές (robust)}
    \item \textgreek{Ο} \textbf{\textgreek{μέσος όρος}} \textgreek{επηρεάζεται πολύ από ακραίες τιμές}
    \item \textgreek{Για δεδομένα με ασυμμετρία, η διάμεσος είναι καταλληλότερη}
\end{itemize}
\end{infobox}

\subsubsection{\textgreek{Γιατί ανά τύπο κρασιού;}}

\textgreek{Τα κόκκινα και λευκά κρασιά έχουν} \textbf{\textgreek{διαφορετικά χημικά χαρακτηριστικά}}:

\begin{table}[H]
\centering
\caption{\textgreek{Σύγκριση διαμέσων τιμών ανά τύπο κρασιού}}
\begin{tabular}{@{}lcc@{}}
\toprule
\textbf{\textgreek{Μεταβλητή}} & \textbf{\textgreek{Κόκκινο}} & \textbf{\textgreek{Λευκό}} \\
\midrule
Fixed acidity & 7.80 & 6.80 \\
Volatile acidity & 0.50 & 0.26 \\
Total sulfur dioxide & ~40 & ~133 \\
Alcohol & 10.20 & 10.40 \\
\bottomrule
\end{tabular}
\end{table}

\textgreek{Αν χρησιμοποιούσαμε τη συνολική διάμεσο, θα εισάγαμε σφάλμα στις εκτιμήσεις.}

\subsubsection{\textgreek{Υλοποίηση}}

\begin{lstlisting}[caption={Fill missing values with median}]
# For each numeric variable
for col in numeric_columns:
    # For each wine type
    for wine_type in ['red', 'white']:
        # Find rows with missing value
        mask = (df['wine_type'] == wine_type) & (df[col].isnull())
        
        # Calculate median for this type
        median_val = df.loc[df['wine_type'] == wine_type, col].median()
        
        # Fill missing with median
        df.loc[mask, col] = median_val

# For wine_type, use mode (most frequent value)
mode_wine = df['wine_type'].mode()[0]  # 'white'
df['wine_type'] = df['wine_type'].fillna(mode_wine)
\end{lstlisting}

% ==============================================================================
\section{\textgreek{Συνοπτικά Αποτελέσματα}}
% ==============================================================================

\subsection{\textgreek{Πίνακας Προβλημάτων και Αντιμετώπισης}}

\begin{table}[H]
\centering
\caption{\textgreek{Συνοπτικός πίνακας προβλημάτων}}
\label{tab:summary}
\begin{tabular}{@{}lcc@{}}
\toprule
\textbf{\textgreek{Τύπος Προβλήματος}} & \textbf{\textgreek{Πλήθος}} & \textbf{\textgreek{Αντιμετώπιση}} \\
\midrule
\textgreek{Περιττές στήλες (Unnamed)} & 2 & \textgreek{Αφαίρεση} \\
\textgreek{Τυπογραφικά λάθη wine\_type} & 9 \textgreek{τύποι (560)} & \textgreek{Διόρθωση/NaN} \\
\textgreek{Μη αριθμητικές τιμές} & 845 & \textgreek{Μετατροπή σε NaN} \\
\textgreek{Αδύνατες/ακραίες τιμές} & 2.037 & \textgreek{Μετατροπή σε NaN} \\
\textgreek{Διπλοεγγραφές} & 1.434 & \textgreek{Αφαίρεση} \\
Missing values (\textgreek{αρχικά}) & 3.092 & \textgreek{Συμπλήρωση με διάμεσο} \\
Missing values (\textgreek{τελικά}) & 0 & -- \\
\bottomrule
\end{tabular}
\end{table}

\subsection{\textgreek{Μεταβολή Μεγέθους Dataset}}

\begin{table}[H]
\centering
\caption{\textgreek{Εξέλιξη μεγέθους dataset}}
\begin{tabular}{|c|c|c|}
\hline
\textbf{\textgreek{Στάδιο}} & \textbf{\textgreek{Γραμμές}} & \textbf{\textgreek{Στήλες}} \\
\hline
\textgreek{Αρχικό dataset} & 8.385 & 15 \\
\hline
\textgreek{Μετά δειγματοληψία (80\%)} & 6.708 & 13 \\
\hline
\textgreek{Μετά αφαίρεση διπλοεγγραφών} & 5.274 & 13 \\
\hline
\textbf{\textgreek{Τελικό καθαρό dataset}} & \textbf{5.274} & \textbf{13} \\
\hline
\end{tabular}
\end{table}

\subsection{\textgreek{Κατανομή Τελικού Dataset}}

\begin{table}[H]
\centering
\caption{\textgreek{Κατανομή τελικού dataset ανά τύπο κρασιού}}
\begin{tabular}{@{}lcc@{}}
\toprule
\textbf{\textgreek{Τύπος Κρασιού}} & \textbf{\textgreek{Πλήθος}} & \textbf{\textgreek{Ποσοστό}} \\
\midrule
\textgreek{Λευκό (white)} & 3.924 & 74.4\% \\
\textgreek{Κόκκινο (red)} & 1.350 & 25.6\% \\
\midrule
\textbf{\textgreek{Σύνολο}} & \textbf{5.274} & \textbf{100\%} \\
\bottomrule
\end{tabular}
\end{table}

% ==============================================================================
\section{\textgreek{Συμπεράσματα}}
% ==============================================================================

\subsection{\textgreek{Τι Μάθαμε}}

\begin{enumerate}
    \item \textgreek{Τα πραγματικά δεδομένα περιέχουν} \textbf{\textgreek{πολλαπλά είδη προβλημάτων}}
    \item \textgreek{Ο καθαρισμός δεδομένων είναι} \textbf{\textgreek{απαραίτητο προκαταρκτικό βήμα}}
    \item \textgreek{Κάθε τύπος προβλήματος απαιτεί} \textbf{\textgreek{διαφορετική αντιμετώπιση}}
    \item \textgreek{Οι αποφάσεις καθαρισμού πρέπει να είναι} \textbf{\textgreek{τεκμηριωμένες και αναπαραγώγιμες}}
\end{enumerate}

\subsection{\textgreek{Περιορισμοί}}

\begin{itemize}
    \item \textgreek{Η συμπλήρωση με διάμεσο} \textbf{\textgreek{μειώνει τη μεταβλητότητα}} \textgreek{των δεδομένων}
    \item \textgreek{Η αφαίρεση διπλοεγγραφών} \textbf{\textgreek{μείωσε το μέγεθος}} \textgreek{του dataset κατά 21\%}
    \item \textgreek{Ορισμένες αποφάσεις βασίστηκαν σε} \textbf{\textgreek{εμπειρική γνώση}}
\end{itemize}

\subsection{\textgreek{Επόμενα Βήματα}}

\textgreek{Το καθαρισμένο dataset είναι πλέον έτοιμο για:}
\begin{itemize}
    \item \textgreek{Περιγραφική στατιστική ανάλυση}
    \item \textgreek{Οπτικοποίηση (boxplots, histograms, κλπ.)}
    \item \textgreek{Εντοπισμό ακραίων τιμών (outliers)}
    \item \textgreek{Κατασκευή προβλεπτικών μοντέλων}
\end{itemize}

% ==============================================================================
\appendix
\section{\textgreek{Παράρτημα: Πλήρης Κώδικας Python}}
% ==============================================================================

\begin{lstlisting}[caption={Complete data cleaning code}]
import pandas as pd
import numpy as np

# Load data
df_original = pd.read_csv('Data Analysis_2026 1st Case_Data.csv')

# Remove extra columns
unnamed_cols = [col for col in df_original.columns if 'Unnamed' in col]
df_original = df_original.drop(columns=unnamed_cols)

# Sampling
AEM = 6931
df = df_original.sample(frac=0.80, random_state=AEM)
df = df.reset_index(drop=True)

# Fix typos in wine_type
wine_type_corrections = {
    'Redd ': 'red', 'r3d': 'red',
    'WHITTE': 'white', 'whit': 'white', 'whate': 'white',
    '999': np.nan, '9999': np.nan, '-999': np.nan, '?': np.nan
}
for wrong, correct in wine_type_corrections.items():
    df.loc[df['wine_type'] == wrong, 'wine_type'] = correct

# Convert to numeric
numeric_columns = ['fixed acidity', 'volatile acidity', 'citric acid',
                   'residual sugar', 'chlorides', 'free sulfur dioxide',
                   'total sulfur dioxide', 'density', 'pH', 'sulphates',
                   'alcohol', 'quality']
for col in numeric_columns:
    df[col] = pd.to_numeric(df[col], errors='coerce')

# Replace impossible values
value_limits = {
    'fixed acidity': (0, 20), 'volatile acidity': (0, 5),
    'citric acid': (0, 3), 'residual sugar': (0, 100),
    'chlorides': (0, 1), 'free sulfur dioxide': (0, 300),
    'total sulfur dioxide': (0, 500), 'density': (0.9, 1.1),
    'pH': (2, 5), 'sulphates': (0, 3), 'alcohol': (5, 20),
    'quality': (0, 10)
}
for col, (min_val, max_val) in value_limits.items():
    mask = (df[col] < min_val) | (df[col] > max_val)
    df.loc[mask, col] = np.nan

# Remove duplicates
df = df.drop_duplicates().reset_index(drop=True)

# Fill missing values
df['wine_type'] = df['wine_type'].fillna(df['wine_type'].mode()[0])
for col in numeric_columns:
    for wine_type in ['red', 'white']:
        mask = (df['wine_type'] == wine_type) & (df[col].isnull())
        median = df.loc[df['wine_type'] == wine_type, col].median()
        df.loc[mask, col] = median

# Save
df.to_csv('cleaned_wine_data_AEM_6931.csv', index=False)
\end{lstlisting}

\end{document}
